\section{Abstract}\label{abstract}

The exponential proliferation of phishing attacks necessitates
proactive, high-precision detection architectures that transcend
traditional reactive blacklisting paradigms. This thesis presents
PhishGuard, a novel, hybrid threat intelligence platform that
orchestrates machine learning (ML), natural language processing (NLP),
and Open-Source Intelligence (OSINT) to mitigate zero-day phishing
campaigns. Operating within an asynchronous Python FastAPI backend, the
system utilizes a heavily optimized XGBoost classifier trained on a
21-dimensional feature vector, achieving a 96.45\% accuracy and a
97.86\% precision rate on a strictly held-out test set of 5,009 samples.
Beyond pure predictive performance, PhishGuard heavily emphasizes
Explainable AI (XAI); it dynamically synthesizes its empirical
classification scores into human-readable heuristic feedback, rendered
through a Next.js frontend, to simultaneously protect and educate the
end-user. This research rigorously details the architectural trade-offs
inherent in synchronous OSINT gathering, the mitigation of concept
drift, and the empirical value of structural versus dynamic intelligence
markers in contemporary cyber defense.

